\documentclass[10pt]{article}
\usepackage[margin=20mm]{geometry}
\usepackage{amsmath,amssymb,mathtools,fontspec,unicode-math,xurl,hyperref}
\IfFontExistsTF{DejaVu Serif}{\setmainfont{DejaVu Serif}}{\setmainfont{Latin Modern Roman}}
\setmathfont{latinmodern-math.otf}
\setlength{\parindent}{0pt}\setlength{\parskip}{6pt}\sloppy
\title{Technical verification of
Maximal commutative unipotent subgroups and a characterization of affine spherical varieties}
\author{AI-assisted technical review}
\date{M25-0278 / focused independent acceptance}
\begin{document}
\maketitle
\section*{Source, scope, and assessment}
The reviewed work is by Andriy Regeta and Immanuel van Santen. The actual source is arXiv:2112.04784v2, 3 August 2023, 30 PDF pages. The frozen DOI is `10.4171/jems/1651`. The EMS landing page records submission on 7 September 2023, acceptance on 17 April 2024, and publication on 8 May 2025. These dates do not turn the earlier author manuscript into the publisher’s version. The arXiv record describes a substantially revised version; it does not label this file an accepted manuscript. Publisher-text equivalence has not been established. All page references below refer to this exact author PDF.\par
Assessment: Major. Theorem A, the principal classification theorem, is false in its stated scope: the trivial unipotent subgroup of a torus satisfies one of its equivalent conditions and fails another. This is a defect of a central conclusion, even though a nontriviality restriction repairs the scope. Two further Major findings concern the stabilizer argument needed by the centralizer proof and the regular-root construction used to synchronize field automorphisms. Four Minor findings have the narrower effects established below. A previously reported eighth item is withdrawn because it misread the source. No counterexample to the spherical reconstruction theorems or Theorem E is claimed.\par
The entire 30-page text was read. Source images on pages 1–28 were actually inspected. The audit concentrated on definitions, closedness, trivial groups, stabilizer families, root dimensions and characters, field twists, the passage to weight monoids, and the dimension restriction in the torus theorem. Reading the whole text is not a claim to have reconstructed every cited theorem.\par
\section*{F1. Major: the central classification includes a false trivial-group case}
Location. Theorem A, p. 2; Proposition 4.1, p. 8 and Claim 1 on p. 9; Corollary 5.2, pp. 12–13. Neither the general conventions nor these statements require the relevant unipotent group to be nontrivial.\par
Work over \(k=\mathbb C\) and take \(X=\mathbb G_m\). Every algebraic \(\mathbb G_a\)-action is trivial. Indeed its pullback sends the unit \(z\) of \(k[z,z^{-1}]\) to a unit of \(k[s,z,z^{-1}]\). Such a unit is \(c z^m\), independent of \(s\). Evaluation at \(s=0\), where the action is the identity, forces \(c=1,m=1\). Thus the action fixes the coordinate ring.\par
Consequently \(G=\{1\}\) is a closed commutative unipotent subgroup and is maximal among commutative unipotent subgroups: there are no other unipotent elements under the definition on p. 2. Condition (b) of Theorem A holds. Condition (a) fails, since \(G\) is properly contained in the commutative group of scalings \(\mathbb G_m\subset\operatorname{Aut}(\mathbb G_m)\). This counterexample satisfies all the global hypotheses, even the stronger normality and uncountability conditions used later. It does not depend on a convention for replicas of the identity.\par
The same omission also affects Proposition 4.1. For \(X=\mathbb A^1,U=\{1\}\), the natural replica construction gives \(R_X(U)=\{1\}\), whereas the centralizer of the identity replicas is all of \(\operatorname{Aut}(\mathbb A^1)\). The exact failure in the proof is visible in Claim 1: it chooses \(u\in U\setminus\{1\}\), a nonzero locally nilpotent derivation \(D_u\), and \(h\) with \(D_u h\ne0\).\par
Counterevidence and repair boundary. For \(U\ne\{1\}\), the cancellation in Claim 1 is valid in the integral coordinate ring, and the rest of the translation argument can be recovered after F2. State Proposition 4.1 for nontrivial \(U\), and Theorem A/Corollary 5.2 for nontrivial \(G\), with nontrivial \(U\) in condition (c). Treat the identity separately: it is maximal unipotent exactly when no nonidentity unipotents exist, and maximal commutative exactly when the whole automorphism group is trivial. The latter follows because the closure of any cyclic subgroup is commutative.\par
This restriction does not by itself refute Corollary 5.3. For a nonidentity unipotent \(u\), its cyclic closure is \(\mathbb G_a\); its replica group is uncountable. Its image has at most countably many connected components, so its identity component is nontrivial. The corrected nontrivial classification is therefore available in that proof; the identity element is handled separately. The Major rating concerns the actual false universal classification, not a claim that every application fails.\par
\section*{F2. Major: a maximal-dimensional stabilizer component need not dominate}
Location. Proposition 4.1, p. 9, the paragraph beginning with the stabilizer family \(E\), after the choice of \(U'\).\par
Consider the faithful action of \(U=\mathbb G_a^4\) on \(X=\mathbb A_x^2\times\mathbb A_t^1\): \[(a,b)\cdot(x,t)=(x+t a+b,t),\qquad a,b,x\in k^2.\] The quotient is \(t\), every fibre is a transitive \(U\)-orbit, and the stated product trivialization is already global over the smooth affine base \(B=\mathbb A^1\). At \(b_0=1,x_0=(0,1)\), the subgroup \(U'=\{(a,0)\}\) complements the stabilizer \(\{(a,-a)\}\). Nevertheless \[E=\{(a,t):t a=0\}
   =(\{0\}\times\mathbb A^1)\ \cup\ (\mathbb A^2\times\{0\}).\] Every fibre is an irreducible linear subspace, and \(E_1=\{0\}\). The component of maximal dimension is the vertical component of dimension two. It does not dominate \(B\), contrary to the assertion used in the proof.\par
Repair and impact. Apply rank openness instead. The infinitesimal action of \(\operatorname{Lie}U'\) along the zero section is injective at \(b_0\), hence on an open neighborhood. In characteristic zero, a zero-dimensional subgroup of a vector group is trivial. On each remaining fibre, commutativity makes stabilizers constant along the transitive \(U\)-orbit; the dimension calculation \(U'+U_x=U\) gives transitivity of \(U'\). The resulting map \(U'\times B\to\pi^{-1}(B)\) is bijective, birational and quasi-finite between smooth varieties, hence an isomorphism by Zariski’s main theorem. The manuscript’s translation argument then continues. The example disproves the intermediate dominance assertion, not the nontrivial centralizer result. The invalid step supplies the free translation model needed to finish Proposition 4.1, on which Theorem A explicitly depends (p. 2). This is a failure in the core derivation, assessed Major even though the rank-openness replacement recovers that step.\par
\section*{Further findings}
\subsection*{F3. Minor: ordinary and generalized root subgroups}
On p. 6 an ordinary root subgroup has dimension one. Proposition 8.1 on p. 18 permits a generalized root subgroup \(V\) of arbitrary positive dimension, but calls its image an ordinary root subgroup of the same dimension. Take \(X=Y=\mathbb A^2\), \(\theta=\mathrm{id}\), \(H=\mathbb G_m\) acting by scalar dilation, and \(V=\mathbb G_a^2\) acting by translations. All hypotheses, including the algebraicity of \(\theta(H)\) and the nontrivial weight, hold. The image has dimension two.\par
The conclusion must say generalized. The proof’s last paragraph on p. 19 establishes the algebraic group and dimension; it supports this reading. To make the scalar-character conclusion explicit, take two coordinate lines in \(V\) and their diagonal. The one-dimensional argument applies to all three. The image of the diagonal is an algebraic \(\mathbb G_a\) subgroup in the product of the two image lines, and both projections are bijective homomorphisms. In characteristic zero they are algebraic isomorphisms. Equivariance forces the two conjugation characters to coincide. Applying this to every pair gives a common nontrivial scalar character on the whole image. This supplies the generalized statement needed later; the example does not disprove it.\par
\subsection*{F4. Major: regular replicas require a nonnegative choice}
In Proposition 6.1(2), pp. 14–15, an integer vector outside finitely many forbidden affine sublattices is chosen and \(f=\prod_j f_j^{a_j}\) is declared regular. Let \(X=\mathbb A^2_{x,y}\), \(H={\mathbb G_a}_x\times{\mathbb G_m}_y\), with translation in \(x\) and dilation in \(y\). This is a connected solvable algebraic group with nontrivial unipotent radical, satisfying the proposition with finite central factor \(\{1\}\). Take \(U={\mathbb G_a}_x,T_1={\mathbb G_m}_y,f_1=y\). The forbidden exponent is zero; \(a_1=-1\) satisfies the printed selection, but \(y^{-1}\notin k[X]\). Its rational replica \((x,y)\mapsto(x+s/y,y)\) does not extend to an automorphism for \(s\ne0\).\par
Also, for higher-dimensional \(T_i\), the forbidden sublattice need only have rank at most \(m-1\). On \(\mathbb A^3_{x,y,z}\) with \(H={\mathbb G_a}_x\times{\mathbb G_m}_y^{}\times{\mathbb G_m}_z^{}\), take \(T_1=T_2=\mathbb G_m^2\), \(f_1=y,f_2=z\), and \(U={\mathbb G_a}_x\). Both forbidden sets are \(\{(0,0)\}\), of rank zero, whereas \(m-1=1\). The statement does not require the \(T_i\) to be distinct.\par
Choose instead a vector in \(\mathbb N^m\) outside the forbidden sets. Each proper affine sublattice is contained in the zero set of a nonzero affine linear polynomial over \(\mathbb Q\). Their product cannot vanish on all of \(\mathbb N^m\), by induction on the number of variables. This simultaneously proves existence, regularity and the required nonzero weights. The proposition’s conclusion survives; the defect is in the stated selection argument. Its use is substantive: on p. 23, Proposition 6.1(2) provides a single root subgroup noncommuting with all the selected tori. Lemma 7.5 then synchronizes the field automorphisms in Theorem 8.5(a), which feeds Theorem 9.1 and Theorem C. Without a regular replica the printed construction does not supply that subgroup. This core construction error is Major; the separate overstatement of the lattice rank is Minor.\par
\subsection*{F5. Minor: exceptional primes for a disconnected diagonalizable group}
The proof of Lemma 8.2, p. 20, identifies the \(p\)-torsion of \(D\) with \((\mathbb Z/p)^{\dim D}\) for every prime \(p\). For \(X=Y=\mathbb A^2,\theta=\mathrm{id}\), the closed diagonalizable subgroup \[D=\{(x,y)\mapsto(t x,\epsilon y):t\in k^*,\ \epsilon\in\{1,-1\}\}
       \simeq\mathbb G_m\times\mu_2\] satisfies every hypothesis. Its dimension is one and its 2-torsion has rank two. Choose a prime not dividing the order of \(D/D^\circ\). The needed torsion-rank bound then holds and still yields \(\dim\theta(D)^\circ\le\dim D\). The conclusion of Lemma 8.2 is not refuted.\par
\subsection*{F6. Minor: two uses of Theorem 8.5 exceed its hypotheses}
In the toric case of Proposition 8.6, p. 25, Theorem 8.5(c) is applied to \(H=DU\). That theorem, p. 22, requires \(H\subseteq G\); here \(G\) is a torus and \(U\simeq\mathbb G_a\) is nontrivial. Already \(X=\mathbb A^1,G=\mathbb G_m,U=\mathbb G_a,D=\{1\}\) realizes all the construction’s hypotheses with \(H\not\subset G\).\par
The same paragraph has already established that \(H\) and \(\theta(H)\) are connected algebraic groups. Proposition 8.1(a), corrected as in F3, therefore applies to \(\theta\) and \(\theta^{-1}\), giving the required dimension-preserving correspondence of generalized roots with nontrivial weight. Roots of trivial weight lie in the respective self-centralizing groups \(D\times\mathbb G_a\); projection to the torus is trivial, so their dimension is at most one. Proposition 8.7 now gives the desired open orbit.\par
In the first case of Theorem E, p. 27, Theorem 8.5(c),(d) is invoked with source \(X=T\) a torus, which 8.5 explicitly excludes. Here \(Y\) is already known to be toric. Instead apply Corollary 5.3 to \(\theta^{-1}\): because \(T\) has no nonidentity unipotents, neither can \(Y\). Lemma 3.4 then forces \(Y\) to be a torus. Both conclusions are retained by these explicit replacements; merely changing a citation without the generalized-root argument would be insufficient.\par
\subsection*{F7. Minor: the dimension statement omits the \(n=1\) boundary}
Theorem 10.1, p. 27, states that a finite-index subgroup \(\Gamma\le\mathrm{SL}_n(\mathbb Z)\) embedding into \(\operatorname{Aut}(Z)\), for irreducible quasi-affine \(Z\), forces \(\dim Z\ge n\). Set \(n=1,\Gamma=\{1\}=\mathrm{SL}_1(\mathbb Z)\), and \(Z=\operatorname{Spec}k\). The index is one, the embedding exists, and \(\dim Z=0<1\).\par
The actual finite-kernel application in Theorem E can only have \(n\ge3\). For \(n=1\) the intermediate integer dimension \(0<r<n\) is impossible. For \(n=2\), \(r=1\) and the target \(\mathrm{GL}_1(\mathbb Z)\) is finite, so the homomorphism from the infinite source cannot have finite kernel. Restricting the auxiliary statement to the used range removes this counterexample. The full external Cantat–Xie theorem, including any \(n=2\) treatment, was not independently reconstructed. No failure of the actual \(n\ge3\) application is asserted.\par
\subsection*{Withdrawn item F8: the postcomposition map was not asserted closed}
The earlier report attributed a closed-immersion assertion to \(\operatorname{End}(W)\to\operatorname{Mor}(W,X)\). The actual text on p. 11 calls this postcomposition map only a morphism. Its closed-immersion assertion concerns instead precomposition \(\operatorname{End}(X)\to\operatorname{Mor}(W,X)\). The proof then obtains a locally closed orbit image and uses that a locally closed subgroup of an algebraic group is closed. The family \(z\mapsto az\) tests the stronger statement that the source did not make. It therefore establishes no defect in Proposition 4.6(2), and F8 is withdrawn.\par
\section*{Positive checks, dependencies, and limits}
For the nontrivial case of Proposition 4.1, fixing all invariant functions forces an automorphism onto the fibres of the generic quotient. After the corrected free \(U'\)-bundle argument, commuting with fibre translations makes it translation by invariant rational coefficients, exactly the prescribed replica product. Such globally regular rational translations are unipotent: on a finite set of coordinate generators, their integer iterates are polynomial in the iterate parameter with rational coefficients. A Vandermonde inversion at distinct integers shows that every coefficient is actually regular. The same polynomial formulas therefore define an algebraic \(\mathbb G_a\)-action. This checks the unipotence assertion used in Remark 4.4.\par
The invariant-field argument of Lemma 5.1 terminates after at most \(\dim X\) strict transcendence-degree drops. The countable-index density step in Corollary 5.3 retains uncountability. The toric lattice calculations retain finite kernel components and use the primitiveness of a root to obtain a connected kernel in Lemma 3.5. The product generation of simple factors and the common nonfinite subgroup in Lemma 7.4 account for compatibility of the field automorphisms in Theorem 8.5; separate factors cannot silently receive unrelated twists.\par
Theorem 9.1 retains the twist \(X^\tau\). Identifying it equivariantly with \(X\) requires the stated \(\mathbb Q\)-defined variety, group and action. Normality is required for the weight-monoid conclusion, and smoothness for the quoted uniqueness result. In Theorem E the hypothesis \(\dim Y\le\dim X\), the finite-index elementary-matrix argument and uncountability all remain essential in the inspected proof.\par
The full Cantat–Regeta–Xie result, the generic-quotient theorem, Borel–Tits, the root-weight/weight-monoid reconstruction and Losev uniqueness, and the Cantat–Xie/Margulis inputs were not independently reproved. A bounded attempt to inspect the external Cantat–Xie PDF returned a publisher 403 and an author-site timeout; these are recorded as failures, not as reading. The projective examples were read but not independently classified in full. No associated separate supplement or computational repository was identified in the recorded searches; this does not establish nonexistence and is not a mathematical defect. No author code was executed.\par
\section*{Supporting checks}
The complete-text reading and original calculations described above belong to the frozen initial review. Its original report, execution records, failures, render-only records and subsequent actual page inspections remain preserved. They do not themselves constitute independent acceptance.\par
The present focused independent acceptance checked all eight previously reported items against this exact source, withdrew F8, and revised the severity and dependency analysis of F2 and F4. The assigned acceptance assistant made these decisions, revised this report, and performed the current compilation, PDF parsing and required actual page inspection. The native bounded read-only assistant `/root/check\_0278\_0279` cross-checked the reported witnesses, hypotheses, counterevidence and source locations. Its exact coverage is recorded separately from the root’s checks in the provenance.\par
Controller UI observations remain attributed to the controller. Runtime instructions identify GPT-6 Astra Pro; backend equivalence was not independently verified. The source version, current report pair, scientific evidence, dependencies, build outputs and page inspection are bound by the current provenance and delivery inventory. Scientific adjudication remains `pending`; `human\_reviews=[]`. No author code was executed.\par
9 A. Regeta and I. van Santen, Maximal commutative unipotent subgroups and a characterization of affine spherical varieties, arXiv:2112.04784v2 (3 August 2023), https://arxiv.org/abs/2112.04784v2. EMS Press, article record for DOI 10.4171/jems/1651, https://ems.press/journals/jems/articles/14298748. This citation supplies publication metadata, not publisher-PDF equivalence.\par
\end{document}
